\documentclass[11pt,a4paper]{article}
\usepackage[margin=1.9cm]{geometry}
\usepackage{fontspec}
\usepackage{enumitem}
\usepackage{xurl}
\usepackage[hidelinks]{hyperref}
\usepackage{xcolor}
\usepackage{titlesec}
\setmainfont{Times New Roman}
\definecolor{heading}{HTML}{003366}
\pagestyle{empty}
\setlength{\parindent}{0pt}
\setlength{\parskip}{5pt}
\sloppy
\setlist[itemize]{leftmargin=1.3em,itemsep=2pt,topsep=2pt}
\titleformat{\section}{\large\bfseries\color{heading}}{}{0pt}{}[\titlerule]
\begin{document}

{\LARGE \textbf{AI / CS Portfolio}}\\
\textbf{Wen Quan}\\
\textcolor{gray}{For PhD Application in Computer Science and Technology, Xiamen University Malaysia}

\section*{Portfolio Positioning}
This portfolio demonstrates my ability to translate built environment research questions into computational, AI-assisted, and web-based assessment systems. My strongest existing demo is the \textbf{China Age-Friendly Community Assessment Map}, which visualizes multidimensional age-friendly community assessment results and provides an interactive evaluation interface.

\section*{Featured Demo}
\textbf{China Age-Friendly Community Assessment Map}

\begin{itemize}
  \item Live demo: \href{https://wenquan0816.github.io/age-friendly-community-map/}{\nolinkurl{wenquan0816.github.io/age-friendly-community-map}}
  \item Repository: \href{https://github.com/WENQUAN0816/age-friendly-community-map}{\nolinkurl{github.com/WENQUAN0816/age-friendly-community-map}}
  \item Type: Web-based interactive assessment dashboard.
  \item Domain: Age-friendly community assessment, spatial visualization, public health, and built environment evaluation.
  \item Status: Published through GitHub Pages and updated to default to English for international review.
\end{itemize}

\section*{Technical Features}
\begin{itemize}
  \item Single-page web application using HTML, CSS, and JavaScript.
  \item Interactive map interface using Leaflet.
  \item Radar chart visualization using ECharts.
  \item Multidimensional evaluation scoring system.
  \item Bilingual interface with English as default.
  \item Community-level data visualization with score categories, filters, and detail panels.
  \item Scoring dimensions include physical space, health services, transport, social participation, safety and emergency, daily services, mental wellness, and governance.
\end{itemize}

\section*{Relevance to Proposed PhD}
The demo is not the final PhD system, but it provides a foundation for the proposed PhD in Computer Science and Technology. It shows that I can:

\begin{itemize}
  \item Build a functional web-based assessment prototype.
  \item Translate age-friendly environment frameworks into structured scoring dimensions.
  \item Present spatial assessment results through interactive visualization.
  \item Communicate complex evaluation outputs to non-technical users.
  \item Develop a research-to-prototype workflow.
\end{itemize}

The proposed PhD will extend this portfolio from community-level assessment toward \textbf{vision-language model assisted indoor space assessment and design recommendation}.

\section*{Proposed Portfolio Extension}
The next version can extend the current demo into:

\begin{enumerate}
  \item \textbf{Indoor Space Risk Assessment Module}: input indoor images, simple room descriptions, or floor-plan screenshots; output age-friendly risk categories and preliminary risk score.
  \item \textbf{Knowledge-Guided Recommendation Module}: use an age-friendly design knowledge base to generate design recommendations such as improving lighting, adding handrails, clearing circulation paths, or modifying furniture layout.
  \item \textbf{Expert Validation Interface}: allow experts to rate whether AI-identified risks and recommendations are correct, useful, and explainable.
  \item \textbf{Comparison Mode}: compare generic VLM output with knowledge-guided output and measure improvements in relevance, consistency, and explanation quality.
\end{enumerate}

\section*{CS/AI Relevance}
This portfolio supports application to Computer Science and Technology because it can be linked to:

\begin{itemize}
  \item Human-computer interaction.
  \item Web-based decision-support systems.
  \item AI-assisted design recommendation.
  \item Spatial data visualization.
  \item Knowledge representation.
  \item Multimodal AI and future VLM integration.
  \item Applied computer vision for built environment assessment.
\end{itemize}

\section*{Selected AI/CS Manuscripts Under Review}
\begin{itemize}
  \item \textbf{Neural Networks}: \textit{Meta-Learning Framework for Healthy Aging Assessment: Attentive Neural Processes with Cross-Region Generalization within CHARLS}. Status: R1 revised submission under processing, NEUNET-D-26-01296R1. This manuscript supports my AI-methods profile through meta-learning, attentive neural processes, and cross-region generalization for healthy aging assessment.
  \item \textbf{Egyptian Informatics Journal}: \textit{Spatial fall risk prediction in elderly residential environments using 3D point cloud-derived features and ensemble learning with SHAP interpretability}. Status: R1 revised submission under processing, EGIJ-D-26-00407R1. This manuscript supports my spatial-intelligence profile through 3D point-cloud features, fall-risk prediction, ensemble learning, and explainable AI.
\end{itemize}

\section*{Application Statement}
I have developed a web-based age-friendly community assessment demo that visualizes multidimensional evaluation results through an interactive map and dashboard. This project demonstrates my ability to translate age-friendly built environment research into computational assessment workflows and user-facing prototypes. For my proposed PhD, I intend to extend this approach toward knowledge-guided vision-language models for indoor space risk assessment and design recommendation.

\end{document}
